\begin{problem}[第32届全国中学生物理竞赛复赛][热力学循环过程]{46}
    如图，1mol 单原子理想气体构成的系统分别经历循环过程 $abcda$ 和 $abc'a$。已知理想气体在任一缓慢变化过程中，压强 $p$ 和体积 $V$ 满足函数关系 $p = f(V)$。
    \begin{subproblem}
        试证明：理想气体在任一缓慢变化过程的摩尔热容可表示为
        \begin{equation*}
            C_{\pi} = C_{V} + \frac{pR}{p + V\frac{dp}{dV}}
        \end{equation*}
        式中，$C_{V}$ 和 $R$ 分别为定容摩尔热容和理想气体常数；
    \end{subproblem}
    \begin{subproblem}
        计算系统经 $bc'$ 直线变化过程中的摩尔热容；
    \end{subproblem}
    \begin{subproblem}
        分别计算系统经 $bc'$ 直线过程中升降温的转折点在 $p-V$ 图中的坐标 $A$ 和吸放热的转折点在 $p-V$ 图中的坐标 $B$；
    \end{subproblem}
    \begin{subproblem}
        定量比较系统在两种循环过程的循环效率。
    \end{subproblem}
\end{problem}